% SOURCE_AUTHORITY: sga5_fr_workpass.tex, lines 7466--7496 (LF-normalized unit).
% SOURCE_SHA256: 791F4EFFC5E02832D5D77ED03518C8156D6F07E4C8238B03545DB93D883FBB28
\begin{prooftext}
A equivalência de (ii) e (iii) é clara, pois \(X\) e \(J^mX\) são estritos. Demonstremos a equivalência de (i) e (ii). Sabe-se (3.2.3) que (i) equivale à existência de um inteiro \(r\) tal que, para todo par \((m,n)\) de inteiros com \(m\ge n+r\), o morfismo de transição que aparece a seguir é um isomorfismo:
\[
X_{m+1}/J^{n+1}X_{m+1}\longrightarrow X_m/J^{n+1}X_m.
\]
Consideremos então o diagrama comutativo exato que aparece a seguir:
\[
\begin{array}{ccccccccccc}
&&0&&0\\
&&\downarrow&&\downarrow\\
&&J^{n+1}X_{m+1}&\longrightarrow&J^{n+1}X_m&\longrightarrow&0\\
&&\downarrow&&\downarrow\\
0&\longrightarrow&F^{m+1}(X_{m+1})&\longrightarrow&X_{m+1}&\longrightarrow&X_m&\longrightarrow&0\\
&&\downarrow&&\downarrow\\
&&X_{m+1}/J^{n+1}X_{m+1}&\longrightarrow&X_m/J^{n+1}X_m&\longrightarrow&0\\
&&\downarrow&&\downarrow\\
&&0&&0 .
\end{array}
\]
Este mostra a existência de um isomorfismo canônico entre
\[
\Ker\bigl(X_{m+1}/J^{n+1}X_{m+1}
       \longrightarrow X_m/J^{n+1}X_m\bigr)
\]
e
\[
F^{m+1}(X_{m+1})/
\bigl(F^{m+1}(X_{m+1})\cap J^{n+1}X_{m+1}\bigr).
\]
Em particular, para que o morfismo de transição após a redução módulo \(J^{n+1}\) seja um isomorfismo, é necessário e suficiente que \(F^{m+1}(X_{m+1})\) esteja contido em \(J^{n+1}X_{m+1}\), o que é precisamente a afirmação (ii).
\end{prooftext}
